\cotefile{29}
\scanpage{127}
\mspage{\handwritten{3}}

De ces deux propriétés résulte formellement qu'on peut trouver un systême projectif \inserted{\handwritten{filtrant}} $(G_i)_{i\handwritten{\in}I}$ dans $\underline{G}$, avec des morphismes $G_i\handwritten{\to}G_j$ qui sont des épimorphismes dans $\unsure{\underline{\mathbb{C}}|\underline{\mathbb{G}}}$, ``essentiellement unique'' dans un sens facile à préciser, tel que \deleted{\unsure{le} foncteur $Z$ soit \unsure{isomorphe}} l'on ait un isomorphisme fonctoriel
\[
Z(G)=\underset{\handwritten{\longrightarrow}}{\lim}\operatorname{Hom}_{\handwritten{\underline{G}}}(G_i,G)
\]

C'est ce systême projectif (pris modulo une équivalence qui \unsure{intuitivement} signifie qu'on a passé à la limite projective des $G_i$) qui peut être noté $\handwritten{\pi}^{\underline{G}}_1(S,a)$ et joue le rôle d'un groupe fondamental de $S$ en $a$ (relativement à la \unsure{catégorìe|catégorie} $\underline{G}$ de groupes, munie du foncteur $F$). Dans le cas \deleted{\unsure{pris}} d'un corps de base $k$, $\underline{G}$ étant la catégorie des groupes algébriques finis sur $k$, on \uline{pourra écrire} $\handwritten{\pi_1}(S/k,a)$,\margin{\handwritten{à tort}} c'est le groupe proalgébrique (mais avec partie infinitésimale) fondamental du $k$-schéma $S$. Lorsque $k$ est parfait, la décomposition d'un groupe algébrique fini en partie infinitésimale et partie réduite montre que le groupe fondamental est produit semi-direct de sa partie réduite ou séparable \deleted{\unsure{définie}} (qui correspond à un groupe compact discontinu ordinaire\handwritten{,} \unsure{surlequel|sur lequel} le groupe fondamental ordinaire de $k$ - i.e. le groupe de Galois de $\bar{k}$ sur $k$ - opère\inserted{\handwritten{\unsure{lorsque|lorsqu'} $a$ est rationnel sur $k$}}), par sa partie infinitésimale, \deleted{\unsure{qui sera un pro-groupe}} qui elle ne dépend d'ailleurs plus du choix du point base $a$ de $S$. Noter que la partie séparable du groupe fondamental peut se réconstituer facilement à l'aide du groupe fondamental ordinaire de $S$ et de son homomorphisme naturel dans $\unsure{\operatorname{gal}|\operatorname{Gal}}(\bar{k}/k)$, \inserted{\handwritten{mais}} il est ``plus grand'' que le groupe fondamental ordinaire, parce\inserted{\handwritten{\,}}qu\deleted{e}\inserted{\handwritten{'}} il correspond à la classification de revêtements principaux de groupe structural non seulement un groupe fini ordinaire, mais un groupe \deleted{\unsure{fini sur lequel}} fini et séparable sur $k$ (i.e. un groupe fini ordinaire sur lequel on fait opérer $\operatorname{Gal}(\bar{k}/k)$ de façon non nécessairement triviale). \deleted{\unsure{Il est}}

J'avoue que si $k$ n'est pas algébriquement clos, le groupe fondamental ci-dessus n'est \deleted{\unsure{peut-être}} \inserted{\handwritten{sans-doute}} pas le bon, on devrait sans doute prendre le groupe fondamental de $S\unsure{\otimes|\times}_{\handwritten{k}}\bar{k}$, qui est un groupe proalgébrique défini sur $\bar{k}$, noter qu'il est muni d'une donnée de descente de $\bar{k}$ à $k$ et est par suite défini en réalité sur $k$; \deleted{avec cette définition, si $S=\operatorname{Spec}(k)$, la \deleted{\unsure{groupe}} partie séparable du groupe fondamental de $S$ (relativement à la ponctuation de $S$ définie par une \deleted{extension \unsure{algé}} clôture algébrique $\bar{k}$ de $k$) n'est autre que}
% archive sheet number: 126
\scanpage{128}
\mspage{\handwritten{4}}

\deleted{$\operatorname{Gal}(\bar{k}/k)$, mais où \unsure{en|on} fait opérer $\operatorname{Gal}(\bar{k}/k)$ par automorphismes intérieurs. Ainsi, \deleted{\unsure{les points}} le groupe des points rationnels sur $k$ du groupe fondamental de $\operatorname{Spec}(k)$ s'identifie au centre de $\operatorname{Gal}(\bar{k}/k)$.}

Voici comment on voit l'existence du systême projectif $(\unsure{\mathbb{E}|\mathbb{B}|E}_i)$. Un couple $(\mathbb{G},z)$, avec $G\in\underline{G}$ et $z\in Z(G)$, est dit \uline{minimal} si on \inserted{ne} peut trouver un sous-groupe $G'\handwritten{\subset}\underline{G}$ de $G$, \deleted{\unsure{tel que}} $\handwritten{\ne}G$, tel que $z$ soit dans de la forme $Z(i)(z')$, avec $z'\handwritten{\in}Z(G')$ et $i:G'\handwritten{\to}G$ l'injection. \deleted{\unsure{Il}} \handwritten{O}n dit qu'un couple $(G,z)$ est \uline{majoré} par un couple $(G',z')$ si on peut trouver un homomorphisme $\inserted{u:}G'\handwritten{\to}G$ tel que $z=Z(u)(z')$. Il résulte de (iii) que tout couple $(G,z)$ est majoré par un couple minimal, et de la propriété 2) de $Z$ que si $(G',z')$ est minimal et domine $(G,z)$, alors il existe \uline{un seul} $u:G'\handwritten{\to}G$ tel que $z=Z(u)(z')$. De ceci résulte que les couples $(G,z)$ minimaux forment un systême projectif pour la relation de domination, d'ailleurs filtrant (car $(G,z)$ et $(G',z')$ sont dominés par $(G\unsure{x|\times}G',(z,z'))$ lui-même dominé par un $(G'',z'')$ minimal), c'est le systême cherché. On peut si on veut choisir un couple $(G,z)$ dans tout systême de couples \inserted{minimaux} isomorphes, de \deleted{\unsure{qui}} telle façon que l'ensemble d'indices \deleted{\unsure{donc}} $I$ devient ordonné et non seulement préordonné. (N.B. Ce genre de raisonnements formels est aussi celui qui sert en théorie des \deleted{\unsure{résidus|réseaux}} modules ...).

Je ne sais encore comment devrait être formulée la suite exacte d'homotopie, pour bien faire l'inclusion dans le groupe fondamental d'une partie infinitésimale devrait donner une théorie satisfaisante du comportement du groupe fondamental par spécialisation. J'espère que si $f:X\handwritten{\to}Y$ est un morphisme propre et séparable à fibres absolument connexes i.e. tel que $f(\underline{O}_X)=\underline{O}_Y$, $X$ étant muni \deleted{\unsure{de}} \inserted{\handwritten{pour simplifier}} d'une section \inserted{\handwritten{$s$}} sur $Y$ qui détermine des points-base sur les fibres, alors les groupes fondamentaux \uline{totaux} des fibres de $X$ forment sur $Y$ un proschéma en groupes \inserted{\handwritten{$\pi_1(X/Y,s)$}} de façon précise qu'on peut trouver un systême projectif \inserted{filtrant} $(G_i)_{i\handwritten{\in}I}$ de schémas en groupes $G_i$ 'spéciaux' sur $\unsure{\mathbb{Y}|\mathbb{S}|\mathbb{X}}$, avec des homomorphismes $G_i\handwritten{\to}G_j$ qui soient des épimorphismes de $\mathbb{Y}$-schémas (i.e. correspondant à un homomorphisme injectif de \unsure{faisceaux} cohérents sur $\mathbb{Y}$), de telle façon que les groupes fondamentaux \unsure{totaux} des fibres de $X$ se\deleted{\unsure{ront}} déduisent dudit pro-schéma en groupes par simple spécialisation. (C'est en tous cas ce que semble indiquer le cas où $X$ est un schéma abélien sur $Y$, cf plus bas).
% archive sheet number: 127
\scanpage{134}
\mspage{I-5}

for $\alpha\handwritten{\ne}o$ the form of the coefficient of $c_\alpha$ implies therefore $\eta_{ij\alpha}=0$ for $i\handwritten{\in}I'$ and $j\handwritten{\in}I''$ since no nonzero linear combination of the $e_j$ is in $W$. Therefore $\xi_{ij}=\eta_{ijo}$ is in $k'$, and $k'$ contains $k(W)$. \hfill q.e.d.

When $V$ is finite dimensional, it follows easily from the previous proof that $k(W)$ is generated by the mutual ratios of the Plücker coordinates of $W$ with respect to any given basis of $V$ rational over $k$.

4.\quad \uline{Rationality for linear and multilinear maps.}

We shall investigate in this paragraph the question of rationality for the multilinear maps. Two particular cases are worth mentioning; the first concerns the linear maps, the second the multilinear and linear forms, that is the case of functions with values in $K$ endowed with the $k$-structure defined by $k$.

Let $V_i$ for $1\leq i\leq r$ and $W$ be vector spaces, each one endowed with a $k$-structure. A multilinear map $f$ from $V_1\times\ldots\times V_r$ to $W$ is called rational over $k$ if $f(v_1,\ldots,v_r)$ is rational over $k$ when the $v_i$ do. It is clear that any multilinear map obtained by composition of some multilinear maps rational over $k$ is rational over $k$.

Let $\{e_\alpha^{(i)}\}$ be a basis of $V_i$ rational over $k$ (for $1\leq i\leq r$); since any vector in $V_i$ rational over $k$ is a linear combination with coefficients in $k$ of the basic elements, a multilinear function $f$ as before is rational over $k$ if and only if $f(e_{\alpha_1}^{(1)},\ldots,e_{\alpha_r}^{(r)})$ is rational over $k$ whatever be the indices $\alpha_1,\ldots,\alpha_r$; if $\{e'_j\}$ is a basis of $W$ rational over $k$ and if one puts
\[
f(e_{\alpha_1}^{(1)},\ldots,e_{\alpha_r}^{(r)}=\sum_j f^j_{\alpha_1,\ldots,\alpha_r}\cdot e'_j
\tag*{(4)}
\]
the function $f$ is rational over $k$ if and only if the components $f^j_{\alpha_1,\ldots,\alpha_r}$ are in $k$. From that follows two facts:

a. Let $k'$ be a field. In order that $f$ be rational over $k'$, it is necessary and sufficient that $k'$ contains the $f^j_{\alpha_1,\ldots,\alpha_r}$, that is the field $k(f)$ generated over $k$ by these elements; in fact, the basis previously introduced are rational over $k$, and \uline{a fortiori} over $k'$.

b. Let us assume the spaces $V_1,\ldots,V_r$ to be finite dimensional. Let $\underline{L}$ be the space of all multilinear maps from $V_1\times\ldots\times V_r$ to $W$, and $\underline{L}_k$ be the

\figure{\begin{tikzpicture}[x=1cm,y=1cm]
\draw[line width=.35pt] (0,.65) -- (.01,.39) -- (0,.17);
\draw[line width=.2pt] (0,.12) -- (0,.08);
\fill (0,0) circle (.012);
\end{tikzpicture}}
% archive sheet number: 133
\scanpage{136}
\mspage{I-4}

$W\handwritten{\cap}V_{k'}=k'\cdot T$ by applying the result a. proved before to $k'$ instead of $k$ and the $k'$-subspace $k'\cdot T$ of $V_{k'}$. Otherwise stated, $W$ is rational over $k'$ and $W\handwritten{\cap}V_{k'}=k'\cdot(W\handwritten{\cap}V_k)$ in such a way that the two $k'$-structures one could get on $W$ by extension of scalars and inducing are actually the same.

We shall now characterize the fields over which a given subspace is rational.

\uline{Proposition} 1.\quad \uline{Let} $V$ \uline{be a vector space with a} $k$-\uline{structure}. \uline{For any subspace} $W$ \uline{of} $V$\unsure{.} \uline{there exists a field} $k(W)$ \uline{such that the relations} ``$W$ is \uline{rational over} $k'$'' \uline{and} ``$k(W)$ \uline{is contained in} $k'$'' \uline{are equivalent for any field} $k'$.

Let $\{e_i\}_{i\handwritten{\in}I}$ be a basis of $V$ rational over $k$; by standard results, there is a subset $I''$ of $I$ such that the vectors $e_j$ for $j$ in $I''$ form a basis of $V$ modulo $W$. This means that $V$ is direct sum of $W$ and the subspace $V''$ having $\{e_j\}_{j\handwritten{\in}I''}$ as a basis; but $V$ is direct sum of $V''$ and the subspace $V'$ having as a basis the $e_i$ for $i$ in $I'=I-I''$. Therefore there exists a unique isomorphism $f$ of $V'$ with $W$ such that $f(v')-v'\handwritten{\in}V''$ for any $v'$ in $V'$. If one puts $v_i=f(e_i)$ for $i\handwritten{\in}I'$, the vectors $v_i$ form a basis of $W$ and they are of the following type
\margin{\figure{\begin{tikzpicture}[x=1cm,y=1cm]
\fill (0,.60) -- (.025,.70) -- (.035,.60) -- (.07,.56) -- (.035,.54) -- (.025,.47) -- (-.015,.48) -- cycle;
\fill (.01,.20) circle (.012);
\draw[line width=.45pt] (.025,.12) -- (.01,.05) -- (-.015,0);
\end{tikzpicture}}}
\[
v_i=e_i-\sum_{j\handwritten{\in}I''}\xi_{ij}\cdot e_j\qquad(i\quad I')
\tag*{(2)}
\]
for some scalars $\xi_{ij}$ in $K$ ($i\handwritten{\in}I'$ and $j\handwritten{\in}I''$). Moreover, no non zero linear combination of the $e_j$ for $j\handwritten{\in}I''$ is in $W$ since $W\handwritten{\cap}V''=0$.

Let $k(W)$ be the field generated over $k$ by the $\xi_{ij}$. If a field $k'$ contains $k(W)$, the vectors $v_i$ are rational over $k'$ and so does $W$. Conversely, suppose that $W$ is rational over some field $k'$. Let $\{c_\alpha\}$ be a basis of $K$ over $k'$ such that $c_o=1$ for some index $\alpha=o$; put $\xi_{ij}=\sum\limits_\alpha\eta_{ij\alpha}\cdot c_\alpha$ with some $\eta_{ij\alpha}$ in $k'$; one gets therefore:
\[
v_i=c_o\cdot(e_i-\sum_j\eta_{ijo}\cdot e_j)+\sum_{\alpha\handwritten{\ne}o}c_\alpha\cdot(\sum_j\eta_{ij\alpha}\cdot e_j).
\tag*{(3)}
\]
Since $W$ is rational over $k'$ and the $v_i$ are in $W$, the considerations in a. above show that the coefficient of each $c_\alpha$ in $v_i$ is in $W$ and rational over $k'$;
% archive sheet number: 135
